\chapter{Data Analysis with jamovi}

A student club surveyed 100 classmates about sleep, screen time, breakfast, and a short quiz. The club wants to know what the answers show. This chapter introduces Jamovi, which is a free software that does the data analysis calculations for you.

\section{Getting Started with jamovi}

\newterm{jamovi}\index{jamovi} is a free statistics program. You choose an analysis from a menu, tick a few boxes, and the results appear right away.

You can use jamovi in two ways:
\begin{itemize}
    \item In a web browser, at \url{https://cloud.jamovi.org}, with a free Guest account. Nothing needs to be installed.
    \item On your own computer, by downloading it from \url{https://www.jamovi.org}. This version works without the internet.
\end{itemize}

The jamovi window has three parts:
\begin{itemize}
    \item On the left, the \newterm{spreadsheet}\index{spreadsheet} holds the data. Each row is one person or thing, and each column is one variable.
    \item Across the top, the \newterm{analysis ribbon}\index{analysis ribbon} lists the analyses you can run.
    \item The \newterm{results panel}\index{results panel} on the right shows tables and graphs.
\end{itemize}

% [figure placeholder: jamovi window with the three parts labeled]

To try it, open an example dataset. Click the menu icon (three lines) in the top-left corner, then choose Open, Data Library, Examples, and Tooth Growth. On the Analyses tab, click Exploration, then Descriptives, and move \texttt{len} into the Variables box. A table of statistics appears without pressing any Run button.

% [C: confirm this menu path in the current jamovi Cloud interface before publication]

\begin{Exercise}[title={Which Part of the Window?}, label={ex:jamovi-regions}]
Name the part of the window you would use for each task: spreadsheet, analysis ribbon, or results panel.

\begin{center}
\begin{tabular}{|p{8cm}|p{4cm}|}
\hline
Task & Part \\
\hline
Fix a number that was typed wrong & \\
\hline
Choose an analysis & \\
\hline
Read the mean & \\
\hline
\end{tabular}
\end{center}
\end{Exercise}
\begin{Answer}[ref={ex:jamovi-regions}]
\begin{center}
\begin{tabular}{|p{8cm}|p{4cm}|}
\hline
Task & Part \\
\hline
Fix a number that was typed wrong & Spreadsheet \\
\hline
Choose an analysis & Analysis ribbon \\
\hline
Read the mean & Results panel \\
\hline
\end{tabular}
\end{center}
\end{Answer}

\section{Setting Up Your Data}

Find \texttt{student\_survey.csv} in your digital resources. Click the menu icon, choose Open, then Browse, and select the file. In jamovi Cloud, you may need to upload it first.

% [C: confirm the Cloud upload and open steps in the current interface ]

\begin{center}
\begin{tabular}{|l|p{9cm}|}
\hline
Column & What it records \\
\hline
\texttt{student\_id} & A code number for each student \\
\hline
\texttt{travel} & How the student gets to school: Walk, Cycle, Bus, or Car \\
\hline
\texttt{breakfast} & Whether the student ate breakfast: Yes or No \\
\hline
\texttt{screen\_time} & Hours on screens per day \\
\hline
\texttt{sleep\_hours} & Hours of sleep per night \\
\hline
\texttt{sleep\_quality} & The student's sleep rating: Poor, Fair, or Good \\
\hline
\texttt{quiz\_score} & Score on a 20-point quiz \\
\hline
\end{tabular}
\end{center}

jamovi gives each column a \newterm{measure type}\index{measure type}, which tells it what kind of data the column holds:
\begin{itemize}
    \item Nominal: categories with no order, such as travel method
    \item Ordinal: categories with an order, such as Poor, Fair, Good
    \item Continuous: measured numbers, such as hours of sleep
    \item ID: labels for each case, such as a student code
\end{itemize}

jamovi guesses these types when it opens a file, and it can guess wrong. To fix one, double-click the column name and pick the correct type.

Two cells in \texttt{sleep\_hours} are empty because those students skipped the question. jamovi treats empty cells as \newterm{missing values}\index{missing values} and leaves them out of calculations.

\begin{Exercise}[title={Measure Types}, label={ex:jamovi-measure-types}]
Give the correct measure type for each column.

\begin{center}
\begin{tabular}{|l|p{4cm}|}
\hline
Column & Measure type \\
\hline
\texttt{student\_id} & \\
\hline
\texttt{breakfast} & \\
\hline
\texttt{sleep\_hours} & \\
\hline
\texttt{sleep\_quality} & \\
\hline
\end{tabular}
\end{center}
\end{Exercise}
\begin{Answer}[ref={ex:jamovi-measure-types}]
\begin{center}
\begin{tabular}{|l|p{4cm}|}
\hline
Column & Measure type \\
\hline
\texttt{student\_id} & ID \\
\hline
\texttt{breakfast} & Nominal \\
\hline
\texttt{sleep\_hours} & Continuous \\
\hline
\texttt{sleep\_quality} & Ordinal \\
\hline
\end{tabular}
\end{center}
\end{Answer}

\section{Describing Data}

Click Exploration, then Descriptives. Move \texttt{quiz\_score} and \texttt{sleep\_hours} into the Variables box.

\begin{center}
\begin{tabular}{|l|r|r|}
\hline
 & \texttt{quiz\_score} & \texttt{sleep\_hours} \\
\hline
N & 100 & 98 \\
\hline
Missing & 0 & 2 \\
\hline
Mean & 12.40 & 7.51 \\
\hline
Median & 12 & 7.60 \\
\hline
Standard deviation & 3.35 & 0.99 \\
\hline
Minimum & 4 & 4.5 \\
\hline
Maximum & 20 & 9.8 \\
\hline
\end{tabular}
\end{center}

The average quiz score was 12.4 out of 20. Students slept 7.51 hours a night on average.

To compare groups, move a category into the Split by box. To draw a graph, open Plots and tick Histogram or Box plot. For a category such as \texttt{travel}, tick Frequency tables to count how many students are in each group.

% [figure placeholder: box plots of quiz_score split by breakfast]

\begin{Exercise}[title={Breakfast and Quiz Scores}, label={ex:jamovi-split}]
Put \texttt{quiz\_score} in the Variables box and \texttt{breakfast} in the Split by box.
\begin{enumerate}
\item What is the median quiz score for each group?
\item Which group did better?
\end{enumerate}
\end{Exercise}
\begin{Answer}[ref={ex:jamovi-split}]
\begin{enumerate}
\item Breakfast: 13 points. No breakfast: 11 points.
\item Students who ate breakfast had the higher median score.
\end{enumerate}
\end{Answer}

\section{Comparing Two Groups}

Students who ate breakfast averaged 13.28 on the quiz, and students who skipped it averaged 11.03. A t-test checks whether a gap this large could be due to chance.

Click T-Tests, then Independent Samples T-Test. Move \texttt{quiz\_score} into Dependent Variables and \texttt{breakfast} into Grouping Variable. Under Tests, tick Welch's and untick Student's. Welch's version is the safer choice because it does not require the two groups to have the same spread.

% [figure placeholder: t-test options panel with Welch's ticked]

jamovi reports $t = -3.40$ and $p = 0.001$. A small $p$-value, below 0.05, means the gap is unlikely to be due to chance alone.

The negative sign appears because jamovi subtracts the Yes group from the No group. The breakfast group still scored higher.

These data come from a survey, so they show a link between breakfast and quiz scores. They cannot show that breakfast causes higher scores.

\begin{Exercise}[title={Reading the t-Test}, label={ex:jamovi-ttest}]
\begin{enumerate}
\item Is $p = 0.001$ below 0.05?
\item Write one sentence explaining what the test shows.
\end{enumerate}
\end{Exercise}
\begin{Answer}[ref={ex:jamovi-ttest}]
\begin{enumerate}
\item Yes.
\item Students who ate breakfast scored higher on the quiz, and the difference is too large to blame on chance alone.
\end{enumerate}
\end{Answer}

\section{Looking at Relationships}

Does more screen time go with less sleep? Click Exploration, then Scatterplot. Put \texttt{screen\_time} on the X-Axis and \texttt{sleep\_hours} on the Y-Axis.
% [C: confirm the Scatterplot item appears under Exploration in the current Cloud version]

% [figure placeholder: scatterplot of sleep_hours against screen_time]

The points slope downward: more screen time goes with less sleep.

Click Regression, then Correlation Matrix, and add both variables. The \newterm{correlation}\index{correlation} $r$ runs from $-1$ to $1$ and measures how closely the points follow a straight line. Here $r = -0.70$, a strong negative relationship.

For a line of best fit, open Regression, then Linear Regression. Put \texttt{sleep\_hours} in Dependent Variable and \texttt{screen\_time} in Covariates. jamovi gives the line
\[
\widehat{\text{sleep}} = 9.55 - 0.44 \times \text{screen time}
\]
Each extra hour of screen time goes with about 0.44 fewer hours of sleep.

\begin{Exercise}[title={Making a Prediction}, label={ex:jamovi-predict}]
Use the line to predict the sleep of a student with 6 hours of screen time.
\end{Exercise}
\begin{Answer}[ref={ex:jamovi-predict}]
$9.55 - 0.44 \times 6 = 6.91$ hours of sleep.
\end{Answer}

\section{Saving Your Work}

Click the menu icon and choose Save As. jamovi saves an \newterm{.omv file}\index{.omv file}, which keeps your data and results together. In jamovi Cloud, download a copy to your own device before closing the tab.
% [C: confirm current Guest plan file-retention behavior in jamovi Cloud before publication]

To put a table or graph in a report, right-click it and choose Copy, then paste it into your document.

\begin{Exercise}[title={Saving Choices}, label={ex:jamovi-saving}]
\begin{enumerate}
\item Which file type keeps both your data and your results?
\item How do you move a table into a report?
\end{enumerate}
\end{Exercise}
\begin{Answer}[ref={ex:jamovi-saving}]
\begin{enumerate}
\item An .omv file.
\item Right-click the table, choose Copy, and paste it into the report.
\end{enumerate}
\end{Answer}


With jamovi you can open data, set it up, describe it, test it, and save your results, all without writing code. Because it is free and runs in a browser, you can use it on almost any device.